\documentclass[a4paper, serif, 10pt]{moderncv}

%% moderncv
% color and style
\moderncvcolor{black}
\moderncvstyle{banking}

%% page layout/margins
\usepackage[top=2.5cm, bottom=2.5cm, left=1.5cm, right=1.5cm]{geometry}

\nopagenumbers{}

%% Babel config for German language
% ref:
% - https://latex3.github.io/babel/guides/locale-german.html
\usepackage[ngerman]{babel}

%% fontspec
\usepackage{fontspec}

\IfFontExistsTF{Linux Libertine}{
  \setmainfont[Ligatures={TeX, Common}, Numbers=OldStyle]{Linux Libertine}
}{}
\IfFontExistsTF{Linux Libertine O}{
  \setmainfont[Ligatures={TeX, Common}, Numbers=OldStyle]{Linux Libertine O}
}{}

%% CTeX
% CJK support, used to show CJK characters in the resume
%
% - heading=false: disable ctex heading
% - punct=kaiming: use kaiming punctuations styles for CJK
% - scheme=plain: use plain scheme, do not override \normalsize font size
% - space=auto: space settings for CJK characters
%
% CJK fonts are resolved per locale language, e.g. Noto Serif/Sans CJK JP
% for Japanese resumes, with fallback to the next region if not installed
%
% ref:
% - http://ctan.mirrorcatalogs.com/language/chinese/ctex/ctex.pdf
\usepackage[UTF8, heading=false, punct=kaiming, scheme=plain, space=auto]{ctex}

\IfFontExistsTF{Noto Serif CJK SC}{
  \setCJKmainfont{Noto Serif CJK SC}
}{}
\IfFontExistsTF{Noto Sans CJK SC}{
  \setCJKsansfont{Noto Sans CJK SC}
}{}

%% line spacing
\usepackage{setspace}
\setstretch{1.125}

%% URL styling - use normal text instead of monospace
\urlstyle{same}

%% Patch \httplink and \httpslink to handle full URLs
% The original moderncv macros always prepend http:// or https:// to URLs.
% This causes issues when the URL already has a protocol (e.g., https://example.com)
% resulting in malformed URLs like https://https://example.com.
% This patch checks if the URL already contains "://" and uses it directly if so.
% Note: We use \str_set:Nx (x-expansion) to expand macros like \@homepage first.
\ExplSyntaxOn
\str_new:N \g_yamlresume_url_str

\RenewDocumentCommand{\httpslink}{O{}m}{
  \str_set:Nx \g_yamlresume_url_str {#2}
  \str_if_in:NnTF \g_yamlresume_url_str {://}
    { \url{#2} }
    { \url{https://#2} }
}

\RenewDocumentCommand{\httplink}{O{}m}{
  \str_set:Nx \g_yamlresume_url_str {#2}
  \str_if_in:NnTF \g_yamlresume_url_str {://}
    { \url{#2} }
    { \url{http://#2} }
}
\ExplSyntaxOff

\name{Anna Schneider}{}
\title{Gesundheits- und Krankenpflegerin mit Schwerpunkt Intensivpflege}
\phone[mobile]{+49 221 9876543}
\email{anna.schneider@pflege-rheinland.de}
\homepage{https://pflege-anna-schneider.de}

\address{Venloer Straße 42, Köln, Nordrhein-Westfalen, Deutschland, 50672}{}{}

\extrainfo{{\small \faLinkedin}\ \href{https://www.linkedin.com/in/anna-schneider-pflege}{@anna-schneider-pflege} {} {} {} • {} {} {} 
{\small \faInstagram}\ \href{https://www.instagram.com/annapflegt}{@annapflegt}}

\begin{document}

\maketitle

\section{Basisinformationen}

\cvline{}{Examierte Gesundheits- und Krankenpflegerin mit über sechs Jahren Erfahrung in der stationären Akut- und Intensivpflege. Ich verbinde fachliche Sicherheit mit einer empathischen, patientenorientierten Haltung und arbeite strukturiert, auch in belastenden Situationen.
%
\begin{itemize}
\item Eigenständige Versorgung und Überwachung intensivpflichtiger Patientinnen und Patienten inklusive Beatmungs- und Medikamentenmanagement
\item Erfahrung in der Anleitung von Auszubildenden sowie in der Umsetzung von Qualitäts- und Hygienestandards
\item Sicher im Umgang mit digitaler Pflegedokumentation und interprofessioneller Zusammenarbeit
\end{itemize}}
\section{Ausbildung}

\cventry{Sept. 2015–Aug. 2019}
        {Bachelor, Pflegewissenschaft, Punkte: 1,7}
        {Hochschule für Gesundheit Bochum}
        {\url{https://www.hsg-bochum.de}}
        {}
        {\begin{itemize}
\item Duales Studium mit integrierten Praxisphasen in Akut- und Langzeitpflege
\item Schwerpunkt auf evidenzbasierter Pflege und Patientenorientierung
\item Bachelorarbeit zum Thema Delirprävention in der Intensivpflege
\end{itemize}
\textbf{Kurse}: Pflegewissenschaftliche Grundlagen, Anatomie und Physiologie, Klinische Pflege, Pflegemanagement, Qualitätsmanagement, Ethik in der Pflege, Gesundheitsförderung}
\section{Berufserfahrung}

\cventry{März 2021–Heute}
        {Gesundheits- und Krankenpflegerin – Intensivstation}
        {Universitätsklinikum Rheinbach}
        {\url{https://www.uniklinikum-rheinbach.de}}
        {}
        {\begin{itemize}
\item Überwachung und Versorgung von Patientinnen und Patienten mit respiratorischer Insuffizienz, Sepsis und postoperativen Komplikationen
\item Bedienung und Überwachung von Beatmungsgeräten, Perfusoren und invasivem Monitoring
\item Anleitung neuer Kolleginnen und Kollegen sowie Auszubildender im Stationsalltag
\item Mitwirkung bei der Weiterentwicklung des Pflegekonzepts und der digitalen Dokumentation
\end{itemize}
\textbf{Schlüsselwörter}: Intensivpflege, Beatmung, Monitoring, Notfallmanagement, Anleitung}

\cventry{Sept. 2019–Feb. 2021}
        {Gesundheits- und Krankenpflegerin}
        {Klinikum St. Elisabeth Köln}
        {\url{https://www.klinikum-st-elisabeth-koeln.de}}
        {}
        {\begin{itemize}
\item Ganzheitliche Pflege und Betreuung internistischer und chirurgischer Patientinnen und Patienten
\item Vorbereitung und Nachsorge bei diagnostischen und therapeutischen Eingriffen
\item Durchführung der Pflegedokumentation und Mitwirkung bei der Pflegevisite
\item Beratung von Angehörigen und Patienten zu Nachsorge und Hilfsmitteln
\end{itemize}
\textbf{Schlüsselwörter}: Akutpflege, Patientenaufklärung, Pflegedokumentation, Angehörigenberatung}
\section{Sprachen}

\cvline{Deutsch}{Muttersprache oder zweisprachig \hfill \textbf{Schlüsselwörter}: Muttersprache}
\cvline{Englisch}{Verhandlungssichere Kenntnisse \hfill \textbf{Schlüsselwörter}: C1, Fachliteratur, Internationale Kommunikation}
\cvline{Spanisch}{Gute Kenntnisse \hfill \textbf{Schlüsselwörter}: Grundkenntnisse}
\section{Fähigkeiten}

\cvline{Intensivpflege}{Experte \hfill \textbf{Schlüsselwörter}: Beatmung, Invasives Monitoring, Sepsisversorgung, Delirprävention}
\cvline{Pflegeprozess}{Erfahren \hfill \textbf{Schlüsselwörter}: Pflegedokumentation, Pflegevisite, Case Management, Entlassmanagement}
\cvline{Notfallmanagement}{Erfahren \hfill \textbf{Schlüsselwörter}: ACLS, Reanimation, Notfallmedikation, Teamkommunikation}
\cvline{Kommunikation}{Erfahren \hfill \textbf{Schlüsselwörter}: Angehörigengespräche, Interprofessionell, Patientenedukation, Konfliktlösung}
\cvline{Qualität und Hygiene}{Fortgeschritten \hfill \textbf{Schlüsselwörter}: Hygienestandards, Qualitätssicherung, Audits, CIRS}
\section{Auszeichnungen}

\cventry{Dez. 2023}
        {Universitätsklinikum Rheinbach}
        {Mitarbeiterin des Jahres 2023}
        {}
        {}
        {Auszeichnung für herausragendes Engagement in der Intensivpflege und in der Anleitung neuer Mitarbeitender.}
\section{Zertifikate}

\cventry{Juni 2023}
        {Deutsche Krankenhausakademie}
        {Fachweiterbildung Intensiv- und Anästhesiepflege}
        {\url{https://www.dka.de/fachweiterbildung}}
        {}
        {}

\cventry{Sept. 2022}
        {European Resuscitation Council}
        {Advanced Cardiac Life Support (ACLS)}
        {\url{https://www.erc.edu/courses/acls}}
        {}
        {}

\cventry{März 2022}
        {Universitätsklinikum Rheinbach}
        {Hygienebeauftragte in der Pflege}
        {}
        {}
        {}
\section{Veröffentlichungen}

\cventry{Mai 2024}
        {Delirprävention in der Intensivpflege – Eine Bestandsaufnahme}
        {Pflegezeitschrift}
        {\url{https://www.pflegezeitschrift.de/artikel/delirpraevention}}
        {}
        {\begin{itemize}
\item Überblick über aktuelle Studien und Leitlinien zur Delirprävention
\item Praxisnahe Empfehlungen für die pflegerische Versorgung auf Intensivstationen
\item Diskussion der Rolle von Angehörigen und interprofessionellen Teams
\end{itemize}}
\section{Projekte}

\cventry{Jan. 2023–Nov. 2023}
        {Projekt zur Umstellung der papierbasierten Pflegedokumentation auf ein digitales System}
        {Einführung der digitalen Pflegedokumentation}
        {}
        {}
        {\begin{itemize}
\item Analyse bestehender Dokumentationsabläufe und Erstellung eines Anforderungskatalogs
\item Schulung von über 40 Pflegekräften im Umgang mit der neuen Software
\item Begleitung der Pilotphase und Auswertung von Anwenderfeedback
\end{itemize}
\textbf{Schlüsselwörter}: Digitalisierung, Pflegedokumentation, Schulung}

\cventry{Apr. 2022–Dez. 2022}
        {Multiprofessionelles Programm zur Früherkennung und Prävention von Delirien}
        {Delirpräventionsprogramm Station 4B}
        {}
        {}
        {\begin{itemize}
\item Einführung eines standardisierten Screenings bei allen über 65-jährigen Patienten
\item Entwicklung von Schulungsmaterialien für Pflege und Medizin
\item Reduktion der Delirinzidenz um 18 Prozent im Projektzeitraum
\end{itemize}
\textbf{Schlüsselwörter}: Delirprävention, Screening, Multiprofessionell}
\section{Interessen}

\cvline{Sport}{Klettern, Laufen, Yoga}
\cvline{Lesen}{Fachliteratur, Romane}
\section{Ehrenamtliche Tätigkeiten}

\cventry{Jan. 2020–Dez. 2023}
        {Ehrenamtliche Sanitätshelferin}
        {Deutsches Rotes Kreuz Köln}
        {\url{https://www.drk-koeln.de}}
        {}
        {\begin{itemize}
\item Absicherung von Veranstaltungen und Unterstützung bei der Ersten Hilfe
\item Betreuung von Verletzten und Koordination der Übergabe an den Rettungsdienst
\end{itemize}}
    

\cventry{Mai 2022–Heute}
        {Mitglied im Ausschuss Pflegequalität}
        {Pflegekammer Nordrhein-Westfalen}
        {}
        {}
        {\begin{itemize}
\item Mitwirkung an Stellungnahmen zu pflegefachlichen Standards
\item Austausch mit Kolleginnen und Kollegen aus Klinik, Pflegeheim und ambulantem Dienst
\end{itemize}}
    
\end{document}